% Corrected French for printed p.214.
% Four print-level repairs are recorded.
\clearpage
\SourcePage{214}

\LayerTarget{tohoku:definition:equivariant-o-homomorphisms}
\noindent l'espace sur lequel on considère $A$ et $B$ (ce qui donne un sens
aussi au symbole $\operatorname{Hom}_{O,G}(U;A,B)$ si $U$ est une partie de
$X$ stable par $G$). On a, en notant que le faisceau
$\mathbf{Hom}_O(A,B)$ est un $G$-faisceau, donc le groupe
$\operatorname{Hom}_O(A,B)=\Gamma_X\mathbf{Hom}_O(A,B)$ un $G$-module, les
formules

\LayerTarget{tohoku:equation:5.6.1}
\[
\tag{5.6.1}
\operatorname{Hom}_{O,G}(A,B)
=\Gamma_X^G\mathbf{Hom}_O(A,B)
=\Gamma^G\operatorname{Hom}_O(A,B).
\]

\LayerTarget{tohoku:definition:equivariant-hom-sheaf}
On pose :

\LayerTarget{tohoku:equation:5.6.2}
\[
\tag{5.6.2}
\mathbf{Hom}_{O,G}(A,B)=(\mathbf{Hom}_O(A,B))^G
\]

\noindent donc $\mathbf{Hom}_{O,G}(A,B)$ est le faisceau sur $Y=X/G$ dont les
sections sur l'ouvert $U\subset Y$ forment le groupe
$\operatorname{Hom}_{O,G}(f^{-1}U;A,B)$ des $G$-$O$-homomorphismes de
$A|f^{-1}(U)$ dans $B|f^{-1}(U)$. Posons :

\LayerTarget{tohoku:definition:four-equivariant-hom-functors}
\LayerTarget{tohoku:equation:5.6.3}
\[
\tag{5.6.3}
\begin{aligned}
\mathbf h_{O,A}(B)&=\mathbf{Hom}_O(A,B),&
h_{O,A}(B)&=\operatorname{Hom}_O(A,B),\\
\mathbf h_{O,G,A}(B)&=\mathbf{Hom}_{O,G}(A,B),&
h_{O,G,A}(B)&=\operatorname{Hom}_{O,G}(A,B),
\end{aligned}
\]

\LayerTarget{tohoku:relation:equivariant-hom-functors}
\noindent définissant ainsi, pour $A$ fixé, quatre foncteurs exacts à gauche,
les deux derniers étant liés aux deux premiers (déjà rencontrés au Chapitre
IV) par les formules fonctorielles

\LayerTarget{tohoku:equation:5.6.4}
\[
\tag{5.6.4}
\mathbf h_{O,G,A}=f_*^G\mathbf h_{O,A}
\]

\LayerTarget{tohoku:equation:5.6.5}
\[
\tag{5.6.5}
h_{O,G,A}=\Gamma_Y\mathbf h_{O,G,A}
=\Gamma^G h_{O,A}=\Gamma_X^G\mathbf h_{O,A}.
\]

\LayerTarget{tohoku:definition:equivariant-ext}
Nous poserons :

\LayerTarget{tohoku:equation:5.6.6}
\[
\tag{5.6.6}
\operatorname{Ext}_{O,G}^n(X;A,B)=R^n h_{O,G,A}(B),
\]

\LayerTarget{tohoku:equation:5.6.7}
\[
\tag{5.6.7}
\mathbf{Ext}_{O,G}^n(A,B)=R^n\mathbf h_{O,G,A}(B).
\]

\LayerTarget{tohoku:property:equivariant-ext-functoriality}
Ainsi, pour deux $G$-$O$-Modules $A,B$, les
$\operatorname{Ext}_{O,G}^n(X;A,B)$ sont des groupes abéliens, formant un
foncteur cohomologique universel en $B$ et se réduisant à
$\operatorname{Hom}_{O,G}(X;A,B)$ en dimension 0, tandis que les
$\mathbf{Ext}_{O,G}^n(A,B)$ sont des faisceaux abéliens sur $Y$, formant un
foncteur cohomologique universel en $B$, et se réduisant à
$\mathbf{Hom}_{O,G}(A,B)$ en dimension 0. D'ailleurs les
$\operatorname{Ext}_{O,G}^n(X;A,B)$ ne sont autres que les groupes $\operatorname{Ext}$
généraux dans la catégorie $C^{O(G)}$. (Bien entendu, nos définitions sont
justifiées grâce à la prop. 5.1.1). Comme $\operatorname{Hom}_{O,G}(A,B)$ donc
aussi $\mathbf{Hom}_{O,G}(A,B)$ est un foncteur contravariant exact en $A$
chaque fois que $B$ est un objet injectif de $C^{O(G)}$, il en résulte
(cf. 2.3) que les $\operatorname{Ext}_{O,G}^n(X;A,B)$ resp. les
$\mathbf{Ext}_{O,G}^n(A,B)$ forment aussi un foncteur cohomologique
contravariant par rapport à $A$ (pour $B$ fixé). Notons enfin qu'on montre
facilement, comme dans 4.2 et 5.2 (qui sont des cas particuliers) que l'on a

\LayerTarget{tohoku:equation:5.6.8}
\[
\tag{5.6.8}
\mathbf{Ext}_{O,G}^n(A,B)|U
=\mathbf{Ext}_{O,G}^n(A|f^{-1}(U),B|f^{-1}(U))
\]

\noindent (\textit{nature locale des} $\mathbf{Ext}$ \textit{par rapport à
l'espace} $Y$). On en conclut un énoncé analogue à celui du corollaire à la
prop. 5.2.1, et qui permet, dans le cas où $G$ est un groupe \textit{discontinu}
d'homomorphismes (cf. 5.3), de se ramener au cas où $G$ est un groupe fini. On
montre alors, en reprenant les raisonnements de 4.1, que les homomorphismes
naturels

\LayerTarget{tohoku:comparison:equivariant-hom-stalk-to-fiber}
\LayerTarget{tohoku:equation:5.6.9}
\[
\tag{5.6.9}
\mathbf{Hom}_{O,G}(A,B)(y)
\longrightarrow \operatorname{Hom}_{U_x}(A(x),B(x))
\qquad (f(x)=y)
\]

\noindent (où $U_x$ est l'anneau, engendré par $O(x)$ et $Z(G_x)$, envisagé
dans 5.1) sont

\TohokuContinuation{215}
